\section{Related Work}

\paragraph{AI-Driven Scientific and Algorithmic Discovery.}
LLM-based discovery systems iteratively generate, evaluate, and refine candidate solutions using prior artifacts and feedback, as in AlphaEvolve~\citep{novikov2025alphaevolve}, OpenEvolve~\citep{openevolve}, CodeEvolve~\citep{assumpccao2025codeevolve}, ShinkaEvolve~\citep{lange2026shinkaevolve}, PACEvolve~\citep{yan2026pacevolve}, DeltaEvolve~\citep{jiang2026deltaevolve} and MLEvolve~\cite{du2026mlevolve}. More recent work emphasizes the importance of exploration itself: SkyDiscover provides adaptive discovery infrastructure~\citep{liu2026skydiscover}, SwarmResearch dynamically orchestrates multiple search branches~\citep{virk2026swarmresearch}, and EvoX~\citep{liu2026evox} explicitly optimizes search strategies rather than only candidate solutions. This shift makes exploration a meta-level optimization problem, but useful supervision for exploration strategies is expensive and delayed because their quality often becomes apparent only after long discovery rollouts. 




\paragraph{Self-Evolving Agents.}
A broader line of work studies agents that improve their own components during interaction. Prior methods evolve model weights~\citep{huang2026r,huang2026g}, agent harnesses~\citep{lee2026meta,zhang2026darwin}, contexts~\citep{zhang2026agentic}, skills~\citep{zhang2026coevoskills,ouyang2026skillos,wu2026agenticrectune}, model behavior through test-time learning~\citep{wang2025thetaevolve,yuksekgonul2026learning,yan2026pacevolve++,wu2026teaching}, rubrics~\citep{xiong2026rubrics}, environments~\citep{huang2026envharness} and other applications~\citep{dai2026takes}. Most operate at the \emph{object level}, improving components used for task execution or reasoning. Recent work has begun to optimize meta-level mechanisms, including search strategies and self-improvement procedures~\citep{liu2026evox,yan2026pacevolve++,wang2026metaskill,zhang2026hyperagents,kim2026meta}. However, such meta-level strategies are difficult to improve because their quality is often revealed only after costly long-horizon rollouts. \textsc{Dream-RSI} makes this meta-level optimization recursive and off-policy by turning accumulated discovery history into replay simulators, allowing exploration controllers to be repeatedly evaluated, improved, and redeployed without rerunning the underlying discovery process.


\paragraph{Memory, History, and Experience Reuse.}
Prior work reuses agent experience as search history, context, memory, reusable skills, or training signals. DeltaEvolve structures evolutionary history through semantic deltas~\citep{jiang2026deltaevolve}; SwarmResearch and MLEvolve use cross-branch or retrospective information to guide subsequent search~\citep{virk2026swarmresearch,du2026mlevolve}; and other work improves how agents access and retain experience through evolving contexts, broader harness state, libraries, or skills~\citep{zhang2026agentic,lee2026meta,xu2026test,ouyang2026skillos}. We take a different view: rather than using exploration history only as context or memory for the next decision, we organize it as a \emph{replay simulator} in which many alternative exploration controllers can be evaluated cheaply. This turns previously collected discovery experience into reusable feedback for meta-level optimization, alleviating the scarcity and high cost of training signals for improving exploration strategies.
